% Created 2026-10-15 Thu 12:50
% Intended LaTeX compiler: pdflatex
\documentclass[presentation]{beamer}
\usepackage[utf8]{inputenc}
\usepackage[T1]{fontenc}
\usepackage{graphicx}
\usepackage{longtable}
\usepackage{wrapfig}
\usepackage{rotating}
\usepackage[normalem]{ulem}
\usepackage{amsmath}
\usepackage{amssymb}
\usepackage{capt-of}
\usepackage{hyperref}
\usepackage[utf8]{inputenc}
\DeclareUnicodeCharacter{2500}{-}
\usepackage{hyperref}
\hypersetup{colorlinks=true, linkcolor=magenta, urlcolor=red}
\RequirePackage{fancyvrb}
\setbeamercolor{titlelike}{fg=white}
\DefineVerbatimEnvironment{verbatim}{Verbatim}{fontsize=\scriptsize}
\usetheme{Warsaw}
\usecolortheme{}
\usefonttheme{}
\useinnertheme{}
\useoutertheme{}
\author{Lenin G. Falconí}
\date{2024-04-27}
\title{Presentaciones Latex con Emacs}
\hypersetup{
 pdfauthor={Lenin G. Falconí},
 pdftitle={Presentaciones Latex con Emacs},
 pdfkeywords={},
 pdfsubject={},
 pdfcreator={Emacs 27.1 (Org mode 9.6.27)}, 
 pdflang={Spanish}}
\begin{document}

\maketitle
\begin{frame}{Outline}
\tableofcontents
\end{frame}



\section{¿Qué es Emacs?}
\label{sec:orgaf099d7}
\begin{frame}[allowframebreaks]{Un entorno de programación extensible y personalizable}
\begin{itemize}
\item No es simplemente un editor de texto.
\item Es un \alert{intérprete de Lisp} diseñado para la computación interactiva.
\item Un entorno de trabajo que evoluciona contigo.
\item "El sistema operativo dentro de tu sistema operativo".
\item Filosofía: Proveer un entorno autocontenido y altamente adaptable
para la manipulación de texto y datos estructurados.
\end{itemize}
\end{frame}

\begin{frame}[allowframebreaks]{Filosofía: Plataforma vs. Producto}
\begin{itemize}
\item \alert{\alert{VS Code:}} Un producto con decisiones predefinidas por terceros.
\item \alert{\alert{Emacs:}} Una plataforma donde tú escribes el comportamiento del software.
\item \alert{\alert{Inversión:}} Tu configuración evoluciona contigo por décadas; no caduca.
\item \alert{\alert{Contexto único:}} "Everything in Emacs" para reducir el cambio de
aplicaciones.
\end{itemize}

Es un ecosistema que puede ser personalizado a tu estilo de trabajo,
con el principio \emph{Everything in Emacs} que reduce la necesidad de
cambiar de contexto entre aplicaciones.
\end{frame}

\begin{frame}[allowframebreaks]{El Corazón: Elisp (Emacs Lisp)}
\begin{itemize}
\item Casi todo lo que ves en Emacs está escrito en Elisp.
\item Puedes evaluar código en tiempo real.
\item No necesitas reiniciar el editor para aplicar cambios profundos.
\item \alert{Ejemplo:} ¿Quieres que el editor se comporte distinto al guardar un archivo? Escribe una función y conéctala a un "hook".
\end{itemize}
\end{frame}

\begin{frame}[fragile,allowframebreaks]{Conceptos Clave de Operación}
 \begin{itemize}
\item \alert{\alert{Búferes y Ventanas:}} La unidad de trabajo es el búfer (archivo, proceso, ayuda).
\item \alert{\alert{Modalidad sin modos estrictos:}} Operación directa complementada con "modos menores".
\item \alert{\alert{Comandos por teclado:}} Uso intensivo de \texttt{M-x} y \texttt{C-x} para ejecutar funciones Elisp.
\item \alert{\alert{Documentación integrada:}} Ayuda instantánea con \texttt{C-h k} (teclas) y \texttt{C-h f} (funciones).
\end{itemize}
\end{frame}

\begin{frame}[fragile,allowframebreaks]{Emacs en el contexto académico y de desarrollo}
 \begin{itemize}
\item \alert{\alert{Redacción técnica:}} Modo \alert{Org} e Integración con \LaTeX y BibTeX (RefTeX).
\item \alert{\alert{Software Libre:}} Garantiza libertad de uso, estudio y
modificación (GNU GPL).
\item \alert{\alert{Gestión bibliográfica:}} Herramientas como \texttt{org-ref} para investigadores.
\item Soporte nativo para decenas de lenguajes de programación.
\end{itemize}
\end{frame}

\begin{frame}[fragile,allowframebreaks]{La Joya de la Corona: Org-mode}
 \begin{itemize}
\item Mucho más que "tomar notas": un ecosistema de productividad completo.
\item \alert{\alert{Gestión:}} Tareas (GTD), agendas y calendarios (\texttt{org-agenda}).
\item \alert{\alert{Exportación:}} Genera PDFs, HTML y diapositivas (como esta) desde un solo origen.
\item \alert{\alert{Literate Programming:}} Documentos que mezclan texto, ecuaciones y código ejecutable.
\end{itemize}
\end{frame}

\begin{frame}[allowframebreaks]{¿Qué más puede hacer Emacs?}
\begin{itemize}
\item \emph{Edición de código}: Completado inteligente (company-mode),
verificación de sintaxis (flycheck), navegación (lsp-mode).
\item \emph{Gestión de proyectos} (projectile, magit).
\item \emph{Composición de documentos} (Org mode, AUCTeX).
\item \emph{Correo y noticias} (mu4e, gnus).
\item \emph{Agenda y calendario} (org-agenda).
\item \emph{Shell integrada} (eshell, term, vterm).
\item \emph{Navegación web} (eww).
\item \emph{Personalización total}: Desde atajos de teclado hasta la
creación de nuevas aplicaciones dentro de Emacs.
\end{itemize}
\end{frame}

\begin{frame}[allowframebreaks]{Conclusión: ¿Qué es Emacs?}
\begin{itemize}
\item Emacs no es solo un editor;
\item Su curva de aprendizaje es pronunciada, pero su retorno en
productividad y control es enorme.
\item \alert{\alert{VS Code:}} Gratificación instantánea, pero techo de personalización bajo.
\item \alert{\alert{Emacs:}} Frustración inicial, pero techo de personalización infinito.
\item Se trata de disfrutar el proceso de "construir tu propia herramienta".
\item Si buscas una herramienta que simplemente funcione de salida, quédate en VS Code.
\item Si buscas una herramienta que sea una \alert{\alert{extensión de tu
pensamiento}} y estás dispuesto a aprender, Emacs es el destino
final.
\end{itemize}
\end{frame}
\section{Objetivos}
\label{sec:org9b4c2cd}
\begin{frame}[label={sec:org8bb6b11},fragile]{Objetivos con Emacs}
 \begin{itemize}
\item En este curso, nos enfocaremos en usar Emacs  para la creación de
presentaciones y artículos académicos con \LaTeX, aprovechando el modo Org
como interfaz unificada. Esto ahorra el uso de la sintaxis de
\LaTeX, reservándola para requerimientos especiales.
\item También abordaremos el concepto de \emph{literate programming} y cómo
\texttt{Org} permite crear documentos que involucran: texto, código y
ecuaciones matemáticas.
\end{itemize}
\end{frame}

\section{Configuración e Instalación en Windows}
\label{sec:org1215567}
\begin{frame}[fragile,allowframebreaks]{Procedimiento de Configuración en Windows}
 \label{org3f7c0a5}

Se requiere realizar los siguientes ajustes para que Emacs opere:

\begin{enumerate}
\item Instalar WSL o un distro de Linux: Se pueden seguir los pasos de la
\href{https://learn.microsoft.com/es-es/windows/wsl/install}{Guía de Instalación}
\item Instalar Emacs:
\begin{verbatim}
        sudo apt-get update
        sudo apt install emacs
\end{verbatim}
\item Instalar \LaTeX{} (básica o completa):
\begin{verbatim}
	sudo apt install texlive-full
\end{verbatim}
Existen otras opciones para Latex que consumen menos recursos (ver
\href{https://linuxconfig.org/how-to-install-latex-on-ubuntu-20-04-focal-fossa-linux}{enlace}):
\begin{verbatim}
        sudo apt install texlive-latex-extra
\end{verbatim}
\end{enumerate}
\end{frame}
\begin{frame}[label={sec:org98938e2},fragile]{Paquetes adicionales de \LaTeX}
 Una vez terminada la instalación principal de \LaTeX instale los
siguientes complementos para no tener problemas en la generación de
archivos PDF desde emacs
\begin{verbatim}
sudo apt-get install texlive-xetex
\end{verbatim}

\begin{verbatim}
sudo apt install texlive-bibtex-extra biber
\end{verbatim}

\begin{verbatim}
sudo apt install latexmk
\end{verbatim}

\begin{verbatim}
sudo apt install texlive-lang-spanish
\end{verbatim}
\begin{enumerate}
\item Instalar Python
\item Configurar el archivo en \textasciitilde{}/.emacs.d/init.el
\end{enumerate}
\end{frame}
\section{Instalación Python}
\label{sec:org487af33}
\begin{frame}[fragile,allowframebreaks]{Instalación de Python}
 Para instalar python se utilizará el entorno de anaconda o
mamba. Existen recomendaciones en favor de mamba. La dirección del
repositorio de Github se encuentra en el enlace \href{https://github.com/conda-forge/miniforge/}{miniforge}

\begin{enumerate}
\item Abra el subsistema  de Linux para Windows: \href{https://learn.microsoft.com/es-es/windows/wsl/install}{guía de instalación}

\item Descargue e instale Mamba: \href{https://github.com/conda-forge/miniforge/releases/latest/download/Miniforge3-Linux-x86\_64.sh}{Descargar ejecutable desde github}
\begin{verbatim}
     wget colocar/direccion/archivo_sh_Linux_x86_64
     sudo chmod +x Miniforge3-Linux-x86_64
     ./Miniforge3-Linux-x86_64
\end{verbatim}
\item Siga los pasos que indica el instalador y reinicie el shell con
\texttt{source \textasciitilde{}/.bashrc} o saliendo de wsl con \texttt{exit} y volviendo a
ingresar desde el \texttt{CMD} de windows con \texttt{wsl -d <Nombre del Distro>}
\item Cree un entorno con las \texttt{python=3.14} para instalar las librerías a utilizar
\begin{verbatim}
     mamba create --name iccd332 python=3.14
\end{verbatim}
\item Para activar el entorno puede usar \texttt{conda activate iccd332} o
\texttt{mamba activate iccd332}. Sin embargo,
\end{enumerate}
\end{frame}
\begin{frame}[fragile,allowframebreaks]{Trouble Shooting Mamba Activate}
 Puede suceder que al ejecutar el comando \texttt{mamba activate} se presente el siguiente error
\begin{verbatim}
  'mamba' is running as a subprocess and can't modify the parent shell.
Thus you must initialize your shell before using activate and deactivate.

To initialize the current bash shell, run:
    $ eval "$(mamba shell hook --shell bash)"
and then activate or deactivate with:
    $ mamba activate
To automatically initialize all future (bash) shells, run:
    $ mamba shell init --shell bash --root-prefix=~/.local/share/mamba
If your shell was already initialized, reinitialize your shell with:
    $ mamba shell reinit --shell bash
Otherwise, this may be an issue. In the meantime you can run commands. See:
    $ mamba run --help

Supported shells are {bash, zsh, csh, posix, xonsh, cmd.exe, powershell, fish, nu}.
critical libmamba Shell not initialized
\end{verbatim}

En este caso se sugiere ejecutar el siguiente comando para poder usar
el comando. Sin embargo, la dirección debe coincidir con la dirección
en donde está instalado el \texttt{miniforge}. Para saber en donde está
instalado el miniforge ejecute:
\begin{verbatim}
conda env list
\end{verbatim}

\begin{verbatim}

# conda environments:
#
# * -> active
# + -> frozen
base                     /home/leningfe/miniforge3
pyt-eqa-fge              /home/leningfe/miniforge3/envs/pyt-eqa-fge
pyt-ml                   /home/leningfe/miniforge3/envs/pyt-ml
tfmlenv              *   /home/leningfe/miniforge3/envs/tfmlenv
whatsapp-osint           /home/leningfe/miniforge3/envs/whatsapp-osint

\end{verbatim}

Para mi caso la dirección es \texttt{/home/leningfe/miniforge3}. Que también
puede escribirse como \texttt{\textasciitilde{}/miniforge3}. En consecuencia, desde el
terminal ejecute:
\begin{verbatim}
  mamba shell init --shell bash --root-prefix=~/miniforge3
\end{verbatim}

Una vez ejecutado reinicie el shell o ejecute \texttt{source
\textasciitilde{}/.bashrc}. Puede reiniciar saliendo de wsl con \texttt{exit} y luego \texttt{wsl
-{}-{}shutdown}, si se abrió sesión de wsl desde un CMD de Windows
\end{frame}
\begin{frame}[label={sec:org3b6ff89},fragile]{Paquetes recomendados de Instalación Python}
 Instalar los siguientes paquetes para disponer del shell avanzado
de Ipython y el complemento \alert{rise} para convertir un jupyter-notebook
en una presentación tipo reveal.js

\begin{verbatim}
mamba install ipython
mamba install jupyter
pip install jupyterlab-rise
\end{verbatim}
\end{frame}
\begin{frame}[label={sec:org0bfbf0d},fragile]{Librerías para Exportar Jupyter Notebook a PDF}
 Para exportar un cuaderno de Jupyter a PDF se requiere que el Sistema
Operativo de Linux disponga de la librería Pandoc:
\begin{verbatim}
sudo apt-get install pandoc
\end{verbatim}
También, si la instalación de \LaTeX se hizo con librerías mínimas, es
necesario disponer de \texttt{xelatex}
\begin{verbatim}
sudo apt-get install texlive-xetex
\end{verbatim}

Con esto la exportación del cuaderno a PDF se realiza haciendo clic en
\alert{\emph{File/Save and Export Notebook As/PDF}}
\end{frame}
\begin{frame}[fragile,allowframebreaks]{Configuración Emacs}
 \label{org454b557}

\begin{itemize}
\item El archivo init.el puede generar instalaciones y activaciones para
emacs por lo que se sugiere reiniciar emacs para que las
configuraciones surtan efecto
\item El comportamiento de Emacs es completamente configurable por el
usuario a través del archivo \textasciitilde{}/.emacs.d/init.el
\item El directorio \texttt{/.emacs.d/} se genera la primera vez que abre
emacs. Debe crear el archivo \texttt{init.el} dentro de éste.
\item Para que no exista conflictos en el uso de la configuración revise
si existe alternativamente el archivo oculto \texttt{.emacs} en la raíz del
sistema operativo. Para esto use el comando \texttt{ls -lha} desde la
posición de \texttt{home}. Por ejemplo, debe existir únicamente el
directorio \texttt{.emacs.d}:
\end{itemize}

\begin{verbatim}
cd
pwd
ls -lha | grep emacs
\end{verbatim}

\begin{itemize}
\item Si existe el archivo \texttt{.emacs} borrarlo con \texttt{rm .emacs} desde la
posición de  home
\end{itemize}
\end{frame}

\begin{frame}[fragile,allowframebreaks]{Configuración Emacs Archivo y creación de archivo init.el}
 \begin{enumerate}
\item Ubíquese en la raíz de Ubuntu \texttt{cd}
\item Verifique que se encuentra en la raíz de su sistema con el comando
\texttt{pwd}. Por ejemplo:
\begin{verbatim}
   cd
   pwd
\end{verbatim}

\item Abra emacs desde la línea de comandos de Ubuntu(WSL)\textasciitilde{}emacs \&\textasciitilde{}

\item Verifique la presencia de archivos ocultos con el comando \texttt{ls -la}
desde la raíz (i.e. \texttt{/home/leningfe})

\item Si existe un archivo \texttt{.emacs} bórrelo usando
\begin{verbatim}
     rm .emacs
\end{verbatim}

\item Verifique que existe únicamente la carpeta \texttt{.emacs.d/}
\begin{verbatim}
   cd 
   pwd
   ls -lha | grep emacs
\end{verbatim}

\begin{verbatim}
/home/leningfe
drwx------ 13 leningfe leningfe 4.0K Oct 15 12:33 .emacs.d
-rw-r--r--  1 leningfe leningfe 157M Apr  9  2025 emacs-config-backup.tar.gz
drwxr-xr-x  2 leningfe leningfe 4.0K Sep 30 15:05 emacs-init-files
\end{verbatim}

\item Abra Emacs con \texttt{emacs \&} desde la raíz

\item Cree el archivo \texttt{init.el} haciendo \texttt{C-x C-f}

\item Escriba el destino en \texttt{/home/<{}<{}User>{}>{}/.emacs.d/init.el}. Sustituya
con \label{orgae61dc5} con su Usuario

\item Copie el contenido del archivo desde \href{https://github.com/LeninGF/EPN-Lectures/blob/main/iccd332ArqComp/configEmacs/init.el}{Github}. Use \texttt{C-y}
para pegar en el archivo creado. Guarde los cambios con \texttt{C-x C-s}.

\item Personalice el archivo con las características personales suyas y
de su máquina.

\item Puede cerrar y reiniciar Emacs para inciar la instalación. Si no
desea salir, puede desencadenar la instalación haciendo \texttt{M-x eval-buffer}
\end{enumerate}
\end{frame}

\section{Comandos básicos de Emacs}
\label{sec:org0554ccf}
\begin{frame}[fragile,allowframebreaks]{Comandos básicos de Emacs}
 Emacs utiliza las teclas \texttt{CTRL} y \texttt{ALT}, en combinación con otras para
desencadenar acciones y comportamientos. Abreviamos \texttt{CTRL} con \texttt{C} y
\texttt{ALT}, que se conoce como \texttt{META}, con la letra \texttt{M}.

Se sugiere que revise el Tutorial Integrado en Emacs para
familiarizarse con el sistema. Este tutorial está disponible en
español ejecutando el comando \texttt{M-x help-with-tutorial-spec-language}.

Para ejecutar un comando compuesto no es necesario mantener presionada
la combinación de teclas. Usted notará en la barra inferior que cuando
ha presionado una combinación, Emacs responde indicando el comando y
aguardando a que complete si es que el comando requiere de
adicionales. Caso contrario se ejecuta.


\begin{center}
\begin{tabular}{|l|l|l|l|}
\hline
Comando & Acción & Comando & Acción \\
\hline
C-x C-s & guardar & C-x C-c & Salir Emacs \\
\hline
C-n & siguiente renglón & C-p & anterior renglón \\
\hline
C-f & avanza un carácter & C-b & atrás un carácter \\
\hline
M-f & avanza una palabra & M-b & atrás una palabra \\
\hline
C-a & ir a inicio & C-e & ir al fin \\
\hline
C-SPC & selecciona región & C-g & cancelar comando \\
\hline
C-d & Borrar un carácter & M-d & borrar palabra \\
\hline
\end{tabular}
\end{center}
\end{frame}
\begin{frame}[label={sec:org85329d3}]{Referencias Web para Comandos Emacs}
\begin{itemize}
\item \href{https://quickref.me/emacs.html}{Emacs Cheat Sheet English}
\item \href{https://cheatography.com/peraza00/cheat-sheets/emacs-guide/\#google\_vignette}{Emacs Cheat Sheet Español}
\item \href{https://www.gnu.org/software/emacs/tour/}{Emacs Tour}
\item \href{https://orgmode.org/manuals.html}{Manual de Documentos Org}
\end{itemize}
\end{frame}
\section{Configuración en Linux}
\label{sec:org6a313fb}
\begin{frame}[label={sec:org8ef5bd1}]{Configuración en Linux}
\begin{description}
\item[{Python}] Se realiza directamente la descarga del Miniforge conforme a lo
señalado en \hyperlink{sec:org487af33}{Instalación Python}
\item[{Emacs}] Se realiza directamente los comandos en  en el punto 3.
\item[{\LaTeX}] Se realiza directamente los comandos en
\hyperlink{org3f7c0a5}{Configuración e Instalación Windows} en el punto 2.
\end{description}

En Linux \alert{\alert{no}} es necesario instalar \alert{\alert{WSL}}
\end{frame}
\begin{frame}[label={sec:org4be19a7},fragile]{Instalación de Corrector Ortográfico}
 Para la corrección ortográfica se requiere de la librería \texttt{aspell}. Esta se puede instalar de la siguiente manera en Linux
\begin{verbatim}
sudo apt install aspell aspell-en aspell-es
\end{verbatim}
\end{frame}
\begin{frame}[fragile,allowframebreaks]{Habilitación de Elpy}
 \begin{itemize}
\item Elpy es un paquete de Emacs para edición de scripts de
Python. Permite completar el código de manera automática, acceder a
documentación, etc.

\item Su auto configuración inicia si creamos un archivo \texttt{.py}. Al hacerlo
pedirá instalar dependencias necesarias. Esta pregunta aparecerá en
la parte inferior en la barra de estado de Emacs y debe contestarla
afirmativamente con una \texttt{y}.

\item Para revisar el estado de la configuración de Elpy ejecute \texttt{M-x
  elpy-config}. Este comando abrirá un buffer en el que indicará el
entorno virtual utilizado, las versiones de python y también si hace
falta instalar alguna paquetería adicional.

\item En el archivo de configuración \texttt{init.el} es importante definir el
\texttt{path} a la carpeta de entornos de \alert{Anaconda}. Para ser cuál es la
carpeta de los entornos ejecute el comando \texttt{mamba info -{}-{}envs}:
\end{itemize}

\begin{verbatim}
mamba env list
\end{verbatim}

\begin{verbatim}
  Name            Active  Path                                         
─────────────────────────────────────────────────────────────────────────
  base                    /home/leningfe/miniforge3                    
  pyt-eqa-fge             /home/leningfe/miniforge3/envs/pyt-eqa-fge   
  pyt-ml                  /home/leningfe/miniforge3/envs/pyt-ml        
  tfmlenv         *       /home/leningfe/miniforge3/envs/tfmlenv       
  whatsapp-osint          /home/leningfe/miniforge3/envs/whatsapp-osint
\end{verbatim}



En el archivo \texttt{init.el} debe personalizar con su dirección respectiva
agregando por ejemplo

\begin{verbatim}
(setenv "WORKON_HOME" "/home/leningfe/miniforge3/envs/")
\end{verbatim}

Reinicie Emacs y luego ejecute \texttt{M-x pyvenv-workon} y seleccione el
entorno de \alert{Anaconda} que se usará para trabajar con Python. Si
ejecuta de nuevo el diagnóstico de Elpy, notará que ya no está con
\texttt{None} en la declaración de \textasciitilde{}Virtualenv\textasciitilde{}/

Instale la paquetería sugerida en el diagnóstico.
\end{frame}

\section{Integración de IA en Emacs GPTEL}
\label{sec:orgd47657e}
\begin{frame}[fragile,allowframebreaks]{Habilitación de Modelos Grandes de Lenguaje o LLM}
 \begin{enumerate}
\item Disponer de una cuenta en GitHub. La configuración del archivo
\texttt{init.el} hace uso por defecto de \texttt{CoPilot}
\item Instalar \texttt{wslview} a fin de que en el primer uso del Chat se pueda
habilitar el uso de \texttt{CoPilot} en el dispositivo: \texttt{sudo apt install wslu}
\item Instalar el paquete para modo de MarkDown. Los archivos de chat con
las \texttt{LLM} se guardan por defecto en formato \texttt{MarkDown}. Sin
embargo, si existiera problema, se pueden modificar para que se
guarden en archivo \texttt{.org}. Para esto puede ejecutar:
\begin{enumerate}
\item \texttt{M-x package-install RET}
\item busca el paquete \texttt{markdown-mode}
\end{enumerate}
\end{enumerate}
\end{frame}
\begin{frame}[fragile,allowframebreaks]{Utilización de Modelos Grandes de Lenguaje o LLM}
 El paquete que realiza la conexión entre \texttt{Emacs} y los diferentes
\texttt{LLMs} es \texttt{GPTEL}. Para iniciar un chat y configurarlo puede ejecutar:
\texttt{M-x gptel-menu}. Se abrirá un menú en donde podrá configurar el
funcionamiento del sistema, cambiar de \texttt{LLM} y configurar desde donde
se envía el prompt y dónde se obtiene la respuesta.

Una sesión de chat se puede grabar \texttt{C-x C-s} y se puede restablecer
abriendo el archivo y activando \texttt{M-x gptel-mode}.
\end{frame}

\section{Configuración e Instalación en MacOS}
\label{sec:org51c5e2c}
\begin{frame}[fragile,allowframebreaks]{Configuración en MacOS (Instalar Brew)}
 \begin{enumerate}
\item \alert{Instalar \texttt{brew}:}
\begin{itemize}
\item Utilizar el comando de la página de HomeBrew
\item El comando identifica la arquitectura del sistema automáticamente
\item Seguir las indicaciones en pantalla
\end{itemize}
\end{enumerate}

\begin{verbatim}
 /bin/bash -c "$(curl -fsSL https://raw.githubusercontent.com/Homebrew/install/HEAD/install.sh)"
\end{verbatim}

\begin{enumerate}
\item Para que \texttt{brew} sea reconocido por el sistema operativo ejecutar:
\end{enumerate}

\begin{verbatim}
 echo >> /Users/delta/.zprofile
 echo 'eval "$(/usr/local/bin/brew shellenv)"' >> /Users/delta/.zprofile
 eval "$(/usr/local/bin/brew shellenv)"
\end{verbatim}

\begin{enumerate}
\item Verificar que \texttt{brew} está funcionado en el sistema usando los
comandos:
\end{enumerate}
\begin{verbatim}
  which brew
\end{verbatim}

\begin{verbatim}
  brew --version
\end{verbatim}
\end{frame}
\begin{frame}[label={sec:org13c2d8d},fragile]{Instalar \LaTeX en MacOS}
 Usando \texttt{brew} proceder a la instalación completa de \LaTeX mediante
el comando:

\begin{verbatim}
  brew install --cask mactex
\end{verbatim}
\end{frame}

\begin{frame}[label={sec:org2b8496b},fragile]{Instalar Emacs en MacOS}
 Utilizar el comando:

\begin{verbatim}
 brew install --cask emacs
\end{verbatim}

Se debe asegurar que Emacs detecte la instalación de \LaTeX
\end{frame}
\begin{frame}[fragile,allowframebreaks]{Instalar Python en Mamba en MacOS}
 \begin{enumerate}
\item Instalar \texttt{wget} o usar e comando \texttt{curl} para descargar el archivo
\end{enumerate}
\begin{verbatim}
  brew install wget
\end{verbatim}
\begin{enumerate}
\item Ejecutar el siguiente comando disponible en el GitHub de Miniforge:
\end{enumerate}
\begin{verbatim}
  wget "https://github.com/conda-forge/miniforge/releases/latest/download/Miniforge3-$(uname)-$(uname -m).sh"
\end{verbatim}
\begin{verbatim}
  curl -L -O "https://github.com/conda-forge/miniforge/releases/latest/download/Miniforge3-$(uname)-$(uname -m).sh"
\end{verbatim}
\begin{enumerate}
\item Verificar la dirección de descarga del archivo
\item Usar \texttt{chmod +x} seguido del nombre del archivo descargado para
convertirlo en ejecutable
\item Realice la ejecución del script de instalación
\item Siga los pasos en pantalla de igual manera  Windows/Linux
\item Reinicie el terminal para observar si \texttt{Mamba} está instalado
\end{enumerate}
\end{frame}

\begin{frame}[fragile,allowframebreaks]{Configuración de Teclado para Emacs}
 Dado que el teclado de MacOs es diferente, es necesario que el archivo
de configuración \texttt{\textasciitilde{}/.emacs.d/init.el} permita el uso adecuado de
\texttt{Command} como \texttt{Alt} o \texttt{Meta} y desactivar el \texttt{Option} a fin de poder
obtener los caracteres especiales del teclado.

\begin{enumerate}
\item Crear o editar el archivo \texttt{\textasciitilde{}/.emacs.d/init.el}
\item Insertar en el archivo la configuración siguiente:
\begin{verbatim}
     (setq mac-command-modifier 'meta)
     (setq mac-option-modifier 'none)
\end{verbatim}
\item Guardar el archivo \texttt{C-x C-s}
\item Evaluar el buffer \texttt{M-x eval-buffer} o reiniciar Emacs
\end{enumerate}
\end{frame}
\begin{frame}[label={sec:orgcdfb0ca},fragile]{Instalar paquetes adicionales}
 Puede usar \texttt{brew} para instalar componentes adicionales como pandoc

\begin{verbatim}
  brew install pandoc
\end{verbatim}
\end{frame}

\section{Presentaciones con Emacs}
\label{sec:orgc5d0ba6}
\begin{frame}[fragile,allowframebreaks]{Hacer una presentación en Emacs}
 \begin{enumerate}
\item Abrir emacs ejecuntando `emacs` desde la línea de comandos
\item Genere un archivo de extensión .org
\begin{itemize}
\item C-x C-f
\item presentacion.org
\end{itemize}
\item Active el modo org para beamer: M-x org-beamer-mode
\item Coloque el encabezado por defecto y el de beamer
\begin{itemize}
\item C-c C-e \# default
\item C-c C-e \# beamer
\alert{\alert{Observar:}} el encabezado de beamer sólo se
activa si el \alert{\alert{modo org-beamer}} está habilitado
\end{itemize}
\item En opciones, cambiar H:1 a H:2 para que los títulos se interpreten
como Secciones de la presentación y las Subsecciones como título de
la presentación
\item Utilice el sistema de marcas de Emacs para escribir la presentación
\item * Sección/Título
\item ** Título Presentación/Subtítulo
\item Una palabra encerrada entre ** equivale a `**negrita**`
\item Una palabra encerrada entre / equivale a `/cursiva/`
\item En el shell, instale el diccionario de español para la corrección ortográfica
\begin{verbatim}
     sudo apt install aspell aspell-es
\end{verbatim}
\item Puede cambiar de diccionario en el buffer usando M-x ispell-change-dictionary
\item Puede activar un modo de revisión al vuelo: M-x flyspell-mode
\end{enumerate}
\end{frame}

\begin{frame}[allowframebreaks]{Algunos Hacks}
\begin{itemize}
\item Para insertar propiedades: C-c C-x p
\item Para insertar código: C-c C-,
\item Para insertar código: M-x org-insert-template-structures y escoge la opción s
\item Para insertar un tag: C-c C-c
\item Para insertar un link: M-x org-insert-link
\item Para generar el archivo pdf: org-beamer-export-to-pdf
\end{itemize}
\end{frame}

\section{Diapositivas de Ejemplo Arq Computadores}
\label{sec:org92b4d5c}
\begin{frame}[label={sec:orgfd7e305}]{El Computador}
Un computador es un sistema complejo conformado por varios subsistemas
electrónicos. Para poder estudiarlo se adopta la perspectiva de
entender su organización y su arquitectura. En general un computador
puede hacer las siguientes operaciones:

\begin{itemize}
\item procesar información
\item almacenar información
\item comunicarse con periféricos de entrada/salida i.e. E/S
\end{itemize}
\end{frame}

\begin{frame}[label={sec:org8e8f4eb}]{El CPU}
La arquitectura principal usada en los computadores es la Von
Neumann. El componente principal del sistema es el CPU, o unidad de
procesamiento central, que es un micro chip. El mismo está formado de:

\begin{itemize}
\item Unidad Aritmética Lógica: ALU
\item Unidad de control
\item Registros
\end{itemize}
\end{frame}

\begin{frame}[label={sec:orgd6093d8}]{A more complex slide}
This slide illustrates the use of Beamer blocks.  The following text,
with its own headline, is displayed in a block:

\begin{theorem}[Org mode increases productivity]
\begin{itemize}
\item org mode means not having to remember \LaTeX{} commands.
\item it is based on ascii text which is inherently portable.
\item Emacs!
\item \(\int e^{-st}f(t)dt\)

\hfill \(\qed\)
\end{itemize}
\end{theorem}
\end{frame}

\begin{frame}[label={sec:orgd343a7f}]{Two Columns}
\begin{columns}
\begin{column}{0.4\columnwidth}
\begin{itemize}
\item this slide consists of two columns
\item the first (left) column has no heading and consists of text
\item second (right) column has an image and is enclosed in an
\alert{example} block
\end{itemize}
\end{column}
\begin{column}{0.6\columnwidth}
\begin{example}[Another block]
\begin{center}
\includegraphics[width=.6\textwidth]{../../images/computer.png}
\end{center}
\end{example}
\end{column}
\end{columns}
\end{frame}

\begin{frame}[label={sec:org601ea05}]{Configurar Org-Babel Python}
Consiste en editar el archivo \textasciitilde{}/.emacs.d/init.el para que emacs tenga
el comportamiento deseado. Para instrucciones referirse al slide \hyperlink{org454b557}{Configuración Emacs}

Este slide muestra cómo hacer una referencia a otra sección
\end{frame}

\begin{frame}[label={sec:org31a1f4d},fragile]{Babel}
 \begin{columns}
\begin{column}{0.45\columnwidth}
\begin{block}{Python code}
\begin{verbatim}
import numpy as np
def greet(name):
    return f"Hello {name}"

return greet("Alice")
\end{verbatim}
\end{block}
\end{column}

\begin{column}{0.4\columnwidth}
\begin{block}{The output}
Hello Alice
\end{block}
\end{column}
\end{columns}
\end{frame}

\section{Funciones avanzadas de WSL}
\label{sec:org3b917ca}

\begin{frame}[label={sec:org622dc5f}]{Funciones avanzadas de WSL}
El uso de espacio de por parte de WSL cambia
dinámicamente. Especialmente, en proyectos de \emph{machine learning}. Para
mitigar esta opción hay las siguientes opciones:

\begin{enumerate}
\item Mover el distro de WSL a otra unidad de disco
\item Comprimir el espacio que ocupa el distro en el disco
\item Eliminar el distro
\end{enumerate}
\end{frame}

\begin{frame}[fragile,allowframebreaks]{Mover el distro de WSL}
 Desde un prompt de Windows realice las siguientes operaciones desde un Powershell
\begin{enumerate}
\item Exportar el distro a otra unidad de disco
\begin{verbatim}
   wsl --export Ubuntu-22.04 D:\WSL\ubuntu22.tar
\end{verbatim}
\item Desinstalar el distro de la unidad \texttt{C:} (Windows)
\begin{verbatim}
   wsl --unregister Ubuntu-22.04
\end{verbatim}
\item Importar de nuevo desde la nueva unidad \texttt{D:}
\begin{verbatim}
   wsl --import Ubuntu-22.04 D:\WSL\Ubuntu22 D:\WSL\ubuntu22.tar --version 2
\end{verbatim}
\end{enumerate}
\end{frame}
\begin{frame}[fragile,allowframebreaks]{Comprimir/Respaldar un distro de WSL}
 \label{org4afddf5}
Realice las siguientes operaciones desde un Powershell
\begin{enumerate}
\item Terminar WSL completamente
\begin{verbatim}
   wsl --shutdown
\end{verbatim}
\item Abrir \texttt{diskpart} en un Powershell
\begin{verbatim}
   diskpart
\end{verbatim}
\item Seleccionar el archivo virtual que contiene el distro. Para esto
debe ubicar la ruta al archivo ext4.vhdx del distro
\begin{verbatim}
    select vdisk file="C:\Users\<YourUser>\AppData\Local\Packages\CanonicalGroupLimited...Ubuntu22.04...\LocalState\ext4.vhdx"
    compact vdisk
\end{verbatim}
\end{enumerate}
\end{frame}
\begin{frame}[label={sec:org8947e26},fragile]{Eliminar el Distro}
 Para eliminar el distro sólo ejecutamos desde un Powershell:
\begin{verbatim}
wsl --unregister <Nombre Distro>
\end{verbatim}

Para esto, el distro debe estar detenido y su nombre localizado. Esto
último lo puede derivar con el comando \texttt{wsl -l -v} desde Powershell
\end{frame}
\begin{frame}[fragile,allowframebreaks]{Actualizar el Distro a una versión nueva: requisitos y respaldo}
 El siguiente procedimiento debe hacerse con cuidado y con un respaldo
de la información ya que puede conllevar la pérdida de datos si no
finaliza con éxito.

El objetivo es actualizar el distro de una version LTS inferior a una
superior (e.g. 22.04 a 24.04).

\begin{enumerate}
\item Antes de comenzar, realizar el procedimiento de
\hyperlink{org4afddf5}{respaldo de la distribución WSL}.

\item En PowerShell de Windows, consultar las distribuciones:

\begin{verbatim}
   wsl --list --verbose
\end{verbatim}

\begin{itemize}
\item Identificar el nombre exacto de la distribución.
\item La columna \texttt{VERSION} debe mostrar \texttt{2}.
\item Si muestra \texttt{1}, detener este procedimiento: primero se debe
convertir la distribución a WSL 2.
\end{itemize}

\item Abrir la terminal de Ubuntu y comprobar su versión:

\begin{verbatim}
   cat /etc/os-release
\end{verbatim}

Para este procedimiento debe indicar \texttt{VERSION\_ID}"22.04"\texttt{.
   Si ya indica =24.04}, no es necesario realizar esta actualización.
\end{enumerate}
\end{frame}

\begin{frame}[fragile,allowframebreaks]{Actualización de Ubuntu: comprobar systemd}
 \texttt{systemd} es el administrador de servicios de Linux. Debe estar
habilitado para seguir el procedimiento de actualización de Ubuntu
en WSL.

\begin{enumerate}
\item Dentro de Ubuntu, ejecutar:

\begin{verbatim}
   ps -p 1 -o comm=
\end{verbatim}

\item Interpretar el resultado:
\begin{itemize}
\item Si devuelve \texttt{systemd}, omitir el apartado de habilitación y
continuar con la actualización.
\item Si devuelve otro nombre, seguir el apartado siguiente.
\end{itemize}
\end{enumerate}
\end{frame}

\begin{frame}[fragile,allowframebreaks]{Actualización de Ubuntu: habilitar systemd si es necesario}
 Realizar este apartado solamente si la comprobación anterior no
devolvió \texttt{systemd}.

\begin{enumerate}
\item En PowerShell, comprobar la versión del programa WSL:

\begin{verbatim}
   wsl --version
\end{verbatim}

El soporte de systemd requiere WSL 0.67.6 o posterior.
Esta versión del programa no es lo mismo que la columna
\texttt{VERSION} de \texttt{wsl -{}-{}list -{}-{}verbose}.

Si WSL es anterior o no reconoce el comando, actualizarlo:

\begin{verbatim}
   wsl --update
\end{verbatim}

\item Dentro de Ubuntu, abrir su archivo de configuración:

\begin{verbatim}
   sudo nano /etc/wsl.conf
\end{verbatim}

\item Añadir lo siguiente, conservando las demás configuraciones:

\begin{verbatim}
   [boot]
   systemd=true
\end{verbatim}

Si ya existe la sección \texttt{[boot]}, añadir o modificar
\texttt{systemd=true} dentro de ella; no duplicar la sección.

\item Guardar y cerrar Nano: \texttt{Ctrl+O}, \texttt{Enter} y \texttt{Ctrl+X}.

\item Guardar cualquier trabajo abierto en otras distribuciones WSL.
En PowerShell, detener WSL:

\begin{verbatim}
   wsl --shutdown
\end{verbatim}

Este comando detiene todas las distribuciones WSL en ejecución.

\item Volver a abrir Ubuntu y comprobar:

\begin{verbatim}
   ps -p 1 -o comm=
\end{verbatim}

Debe devolver \texttt{systemd}. Si no lo hace, no continuar con la
actualización hasta resolver este requisito.
\end{enumerate}
\end{frame}

\begin{frame}[fragile,allowframebreaks]{Actualización de Ubuntu: ejecutar y verificar}
 \begin{enumerate}
\item Dentro de Ubuntu, actualizar los paquetes de la versión actual:

\begin{verbatim}
   sudo apt update && sudo apt full-upgrade -y
\end{verbatim}

Esperar a que termine correctamente. Si presenta errores,
resolverlos antes de continuar.

\item Iniciar la actualización de versión:

\begin{verbatim}
   sudo do-release-upgrade
\end{verbatim}

Confirmar que el destino propuesto sea Ubuntu 24.04 y seguir
las instrucciones. No interrumpir el proceso.

\item Al finalizar, seguir las instrucciones de reinicio del
actualizador. Si es necesario reiniciar la instancia desde
PowerShell, usar:

\begin{verbatim}
   wsl --terminate "NOMBRE_DE_LA_DISTRIBUCION"
\end{verbatim}

Sustituir \texttt{NOMBRE\_DE\_LA\_DISTRIBUCION} por el nombre obtenido
con \texttt{wsl -{}-{}list -{}-{}verbose}. Después, volver a abrir Ubuntu.

\item Verificar el resultado dentro de Ubuntu:

\begin{verbatim}
   cat /etc/os-release
   ps -p 1 -o comm=
\end{verbatim}

Los resultados esperados son:
\begin{itemize}
\item \texttt{VERSION\_ID}"24.04"=.
\item \texttt{systemd} como proceso principal.
\end{itemize}

El nombre registrado de la distribución en WSL puede seguir
siendo el anterior; comprobar la versión con \texttt{os-release}.

\item Probar las herramientas habituales: Emacs, exportación de
documentos Org a PDF y los entornos de Python.
\end{enumerate}

\begin{block}{Documentación}
\begin{itemize}
\item \href{https://ubuntu.com/wsl/docs/stable/howto/upgrade-ubuntu/}{Ubuntu: actualizar Ubuntu en WSL}
\item \href{https://learn.microsoft.com/en-us/windows/wsl/systemd}{Microsoft: habilitar systemd en WSL}
\end{itemize}
\end{block}
\end{frame}

\section{Funciones avanzadas de Anaconda}
\label{sec:org73ebe5a}
Los entornos de anaconda crecen con la instalación de paquetes y
librerías. A fin de evitar descargar nuevamente de internet, anaconda
almacena las descargas en el disco. Esto ocupa demasiado espacio, por
lo que es recomendable limpiar el cache, los logs y paquetes no
usados.
\begin{frame}[label={sec:org9a68bb4},fragile]{Comandos de Limpieza para Conda}
 Ejecute el siguiente comando para eliminar \texttt{tarballs}, limpiar el
cache, logs y paquetes no usados.
\begin{columns}
\begin{column}{0.4\columnwidth}
\begin{verbatim}
conda clean --all --yes
\end{verbatim}
\end{column}
\begin{column}{0.6\columnwidth}
\begin{center}
\includegraphics[width=.9\linewidth]{/mnt/d/entea/Lectures/iccd332ArqComp/images/conda-clean.png}
\end{center}
\end{column}
\end{columns}
\end{frame}
\begin{frame}[label={sec:org69436cc},fragile]{Eliminar un entorno de Conda}
 Para eliminar el entorno es necesario primero salir de él con el
comando \texttt{conda deactivate}. Una vez que esta desactivado:
\begin{verbatim}
conda deactivate
conda remove --name myenv --all
\end{verbatim}
Reemplace \texttt{myenv} con el nombre del entorno. Esto borrará toda la
información del entorno.
\end{frame}
\begin{frame}[label={sec:org8f080ce},fragile]{Crear un backup del entorno}
 Para replicar resultados es importante trabajar con los mismos
paquetes (incluyendo las versiones específicas de los mismos). Esto se
puede lograr exportando a un archivo \texttt{YAML}:
\begin{verbatim}
conda activate myenv
conda env export > myenv.yml
\end{verbatim}

Note que las versiones de los paquetes cambian con el sistema
operativo por lo que si el proyecto trabaja en un S.O. Windows, las
versiones no serán las mismas que se instalan en un S.O. Linux. Una
alternativa para evitar este problema y garantizar replicación es
trabajar con \texttt{Docker}.

Verifique que el archivo contenga tanto los paquetes de \texttt{pip} como los
de \texttt{anaconda}. 
\end{frame}
\begin{frame}[label={sec:orgbc837b1},fragile]{Crear un entorno desde YAML}
 \begin{verbatim}
conda env create -f myenv.yml
\end{verbatim}
Para verificar los paquetes instalados o restituidos use
\begin{verbatim}
conda list
\end{verbatim}
\end{frame}
\begin{frame}[fragile,allowframebreaks]{Entorno Anaconda Portable}
 El entorno se puede exportar como archivo portable para moverlo a otro
sistema operativo (de igual naturaleza) con:
\begin{verbatim}
conda pack -n myenv -o myenv.tar.gz
\end{verbatim}
Para restituirlo:
\begin{verbatim}
mkdir -p ~/envs/myenv
tar -xzf myenv.tar.gz -C ~/envs/myenv
source ~/envs/myenv/bin/activate
\end{verbatim}
\alert{\alert{Importante}} debe estar instalado en el entorno el paquete \texttt{conda
  pack}. Se puede instalar de las siguientes maneras:
\begin{verbatim}
    conda install -c conda-forge conda-pack
    mamba install conda-pack
    pip install conda-pack
\end{verbatim}
\end{frame}
\end{document}
